% latest version : https://www.overleaf.com/project/6a2fd0023e497c7e3ba3b0cb 

\documentclass[11pt, a4paper]{article}

\usepackage{amsmath}
\usepackage{amssymb}
\usepackage{graphicx}
\usepackage{booktabs}
\usepackage{hyperref}
% Margins
\usepackage[top=2.5cm, bottom=2.5cm, left=2.5cm, right=2.5cm]{geometry}

% New commands for tool names
\newcommand{\ape}{\textit{ape}}
\newcommand{\dendropy}{\textit{DendroPy}}
\newcommand{\cogent}{\textit{cogent3}}
\newcommand{\biopython}{\textit{Biopython}}
\newcommand{\ete}{\textit{ete3}}
\newcommand{\toytree}{\textit{toytree}}
\newcommand{\phylo}{\textit{Phylo.jl}}
\newcommand{\gotree}{\textit{gotree}}
\newcommand{\phylors}{\textit{phylo-rs}}
\newcommand{\castor}{\textit{castor}}
\newcommand{\stan}{\textit{Stan}}
\newcommand{\turing}{\textit{Turing.jl}}
\newcommand{\biomc}{\textit{BIOMEC}}

\title{A Comparative Analysis of Phylogenetic Tree Manipulation Tools: \\
Benchmarking Tree Traversal and Coalescent Likelihood}

\author{Emma Gairaud Hess}

\date{\today}

\begin{document}

\maketitle

% ==========================================
% ABSTRACT
% ==========================================
\begin{abstract}

Phylogenetic tree manipulation is essential for evolutionary biology and epidemiology.
This study presents a survey and benchmarking of 8 phylogenetic packages and tools,
evaluating their performance on tree traversal operations (pre-order, post-order, level-order) and ultrametricity testing.
The tests were conducted on randomly generated trees of 1{,}000 leaves.
Performance metrics were measured in R-based (\ape, \castor), Python-based
(\dendropy, \cogent, \biopython, \ete, \toytree) and Julia-based (\phylo) tools.
Subsequently, the best-performing tools were selected, and a coalescent likelihood framework was constructed around them, deriving the full log-likelihood of Kingman's n-coalescent for both constant and variable effective population sizes. 
This framework was then extended to support Bayesian inference via Hamiltonian Monte Carlo (HMC) using probabilistic programming languages (\stan, \turing), with applications in epidemic outbreak modeling and parameter estimation.

! leaving out BIOMEC for now since it seems to be specific to electrochem (potential replacement: pyhmc)
\end{abstract}

\newpage
\tableofcontents
\newpage

% ==========================================
% 1. INTRODUCTION
% ==========================================
\section{Introduction}

Phylogenetic trees are fundamental data structures in evolutionary biology, representing the evolutionary relationships between taxa. Modern phylogenetic studies frequently involve large trees (thousands to millions of leaves), whose manipulation requires efficient algorithms.

! cite covid or other article containing >1mill samples (ie https://pmc.ncbi.nlm.nih.gov/articles/PMC10131274/)

In addition, coalescent processes are used in epidemiology to provide a framework for estimating parameters such as transmission rates and the number of unsampled individuals during disease outbreaks \cite{volz2009,pybus2001}. These models require extensive tree traversal and statistical inference, making computational efficiency crucial as well.

Current phylogenetic software is highly diverse, with implementations available across multiple programming languages. While many studies focus on phylogenetic inference accuracy, comparative performance evaluation of fundamental tree operations is limited.

!! read up PhyloSmew (benchmark BUT for accuracy not computational cost/efficiency)

This report addresses two key objectives:
\begin{enumerate}
    \item Benchmarking basic tree operations across several phylogenetic packages.
    \item Establishing a mathematical and computational framework for coalescent likelihood evaluation. 
\end{enumerate}

% ==========================================
% 2. METHODS
% ==========================================
\section{Methods}

% ------------------------------------------
\subsection{Part I: Tree Traversal Benchmarking}
% ------------------------------------------

\subsubsection{Benchmark Design}

A standardized benchmarking framework was implemented to evaluate four core operations:
\begin{itemize}
    \item \textbf{Preorder traversal:} Root-to-leaf depth-first search (DFS) (fig.~1A)
    \item \textbf{Postorder traversal:} Leaf-to-root DFS (fig.~1B)
    \item \textbf{Level-order traversal:} Queue-based breadth-first search (BFS) (fig.~1C)
    \item \textbf{Ultrametricity test:} Verifies that all leaves are equidistant
          from the root (fig.~2)
\end{itemize}

! insert R figures from searchillust.r
! FIGURE 1: preorder (A), postorder (B) and level order (C) traversal on the same tree
! FIGURE 2: ultrametric (left) vs heterochronic (right) tree

\subsubsection{Tree Generation Procedure}

Test trees were generated in \ape, with a set seed, by using coalescent models of 1{,}000 leaves:
\begin{itemize}
    \item \textbf{Ultrametric trees:} \texttt{rcoal(n)} generates a coalescent tree of $n$ leaves by randomly clustering its tips
    \item \textbf{Heterochronic trees:} \texttt{rtree(n)} generates a coalescent tree of $n$ leaves by randomly splitting its edges
\end{itemize}

! find clearer definitions for both

For each tree type, 1{,}000 trees were generated and stored in \texttt{.nwk} format files, guaranteeing the use of an equal tree set across all tested tools. 

\subsubsection{Benchmarking Procedure}

Timing measurements were done after a warmup iteration to account for JIT compilation and memory allocation overhead. Execution times were recorded in microseconds using
\texttt{Sys.time()} (R), \texttt{time.perf\_counter()} (Python) or \texttt{time\_ns()} (Julia).

! despite the warmup there are probably some remaining initialization issues, need to test with a different tree set and/or modify warmup logic

! mention that some time functions are not in microseconds, so they were converted manually (ie time * $10^6$ s -> micros)?

\subsubsection{Tools Evaluated}

\paragraph{\ape} (R package)

An R package for phylogenetic analysis. It provides native tree data structures using edge lists stored in $n \times 2$ matrices. Tree operations are implemented in compiled C code for performance. Benchmarks use the \texttt{is.ultrametric()} function from the compiled backend.

\paragraph{\castor} (R package)

Specialized R package for large-scale tree inference and analysis. \castor\ stores trees using the same \ape-compatible \texttt{phylo} objects, while adding additional routines for ancestral state reconstruction and diversification modeling.

\paragraph{\dendropy} (Python)

Python implementation of phylogenetic data structures. \dendropy\ uses object-oriented design with \texttt{Node} and \texttt{Tree} classes. It has no compiled backend, so performance is bottle-necked by the Python interpreter overhead.

\paragraph{\cogent} (Python)

It's designed for comparative genomics with phylogenetic capabilities. Tree operations include optional Cython compilation for performance-critical sections and NumPy-based array operations where applicable.


\paragraph{\biopython} (Python)

A general-purpose bioinformatics package with phylogenetic submodule (\texttt{Bio.Phylo}). Provides wrappers around external phylogenetic tools while offering native Python tree structures. Emphasises compatibility over computational efficiency.

\paragraph{\ete} (Python)

Python implementation optimised for interactive phylogenetic analysis. Recent versions include compiled extensions for critical operations. The library emphasises usability and visualisation alongside performance.

\paragraph{\toytree} (Python)

Lightweight Python tree library with Cython bindings, providing significant speedup over pure Python alternatives. Can be integrated with Jupyter notebooks for interactive workflows.

\paragraph{\phylo} (Julia)

A Julia package providing multiple-dispatch tree implementations. Julia's JIT compilation enables performance comparable to compiled languages.

! this is where the post-order max is extremely high; TO DO: add another nwk file with different trees to see if issue persists (added max index print -   Max at index:  987; likely not a JIT/compiler issue then)

\subsubsection{Statistical Analysis}

For each tool and operation, summary statistics were computed:
\begin{itemize}
    \item Mean, median, and standard deviation of execution times
    \item Min/max times to identify outliers and JIT compilation effects
\end{itemize}

Summary visualisations (violin plots) were generated for qualitative comparison.

Tools presenting outlier max values had a max index tracker to verify the source of the problem.

% ------------------------------------------
\subsection{Part II: Coalescent Likelihood Framework}
% ------------------------------------------

\subsubsection{Biological Motivation}

Kingman's $n$-coalescent \cite{kingman1982} is a stochastic process that describes the genealogical ancestry of a random sample of $n$ individuals drawn from a large population assuming exchangeable reproduction, going backwards in time. It is the mathematical foundation for likelihood-based inference of effective population size $N_e$ from phylogenetic trees, and is closely related to birth-death models used in epidemiology and other fields.

! review that exchangeability in kingman's paper does refer to the populations reproduction (and if it's "neutral" evolution or something else) 
\subsubsection{Kingman's Coalescent Process}

Consider a sample of $n$ lineages observed at time $t = 0$ and going backwards in time. Under the standard Wright-Fisher model with effective population size $N_e$, the probability that any specific pair of lineages coalesces in the previous generation
is $1/N_e$. The waiting time until the first coalescence event among $k$ extant lineages follows an
exponential distribution.

\medskip
\noindent\textbf{Coalescence rate.}
When $k$ lineages are present, there are $\binom{k}{2} = k(k-1)/2$ distinct pairs,
each coalescing at rate $1/N_e$, giving a total instantaneous coalescence rate of: 

\begin{equation}
    \gamma_k = \frac{\binom{k}{2}}{N_e} = \frac{k(k-1)}{2N_e}.
    \label{eq:coalrate}
\end{equation}

\noindent\textbf{Waiting time distribution.}
Denoting the waiting time during which exactly $k$ lineages are present as $w_k$,
we have
\begin{equation}
    w_k \sim \mathrm{Exp}\!\left(\gamma_k\right),
    \qquad
    f(w_k) = \gamma_k\, e^{-\gamma_k w_k},
    \quad w_k \geq 0.
    \label{eq:waitingtime}
\end{equation}

Starting with $n$ lineages, the process generates $n - 1$ coalescent events, reducing
the number of lineages from $n$ down to $1$ (the most recent common ancestor, MRCA).
The topology of which two lineages merge at each event is chosen uniformly at random
from the $\binom{k}{2}$ possible pairs.

\subsubsection{Coalescent Likelihood Under Constant Population Size}

Let $\mathcal{T}$ denote a fully resolved rooted binary tree with $n$ leaves and
coalescent waiting times $w_n, w_{n-1}, \ldots, w_2$, where $w_k$ is the duration
during which $k$ lineages coexist (measured backwards from the sampling time).
Assuming a constant effective population size $N_e$, the likelihood of the tree
given $N_e$ factorises over coalescent intervals:

\begin{equation}
    \mathcal{L}(N_e \mid \mathcal{T})
    = \prod_{k=2}^{n}
      \frac{\binom{k}{2}}{N_e}\,
      \exp\!\left(-\frac{\binom{k}{2}}{N_e}\,w_k\right).
    \label{eq:coallikelihood}
\end{equation}

Taking logarithms yields the log-likelihood

\begin{equation}
    \ell(N_e \mid \mathcal{T})
    = \sum_{k=2}^{n}\left[\log\binom{k}{2} - \log N_e
      - \frac{\binom{k}{2}}{N_e}\,w_k\right]
    = (n-1)\bigl(-\log N_e\bigr) + C - \frac{S}{N_e},
    \label{eq:loglik}
\end{equation}

where $C = \sum_{k=2}^{n}\log\binom{k}{2}$ is a constant independent of $N_e$, and

\begin{equation}
    S = \sum_{k=2}^{n} \binom{k}{2}\,w_k
    \label{eq:sstat}
\end{equation}

is the sufficient statistic for $N_e$: the total coalescent-weighted tree length.

\medskip
\noindent\textbf{Maximum likelihood estimator.}
Setting $\partial\ell/\partial N_e = 0$ gives the closed-form MLE

\begin{equation}
    \hat{N}_e = \frac{S}{n-1}
              = \frac{1}{n-1}\sum_{k=2}^{n}\binom{k}{2}\,w_k.
    \label{eq:MLE}
\end{equation}

This estimator has an intuitive interpretation: $\hat{N}_e$ is the average
coalescent-pair--weighted waiting time per coalescent event.

\subsubsection{Extension to Variable Effective Population Size}

In practice, $N_e$ may change over time, for example due to population expansions or
epidemic dynamics. Griffiths and Tavaré (1994) \cite{griffiths1994} generalised the
coalescent to allow a time-varying effective population size $N_e(t)$, where $t$
increases into the past from the sampling time $t = 0$.

\medskip
\noindent\textbf{Likelihood with $N_e(t)$.}
When $k$ lineages are present starting at backward time $t_k$, the probability density
of the waiting time $w_k$ is

\begin{equation}
    f\!\left(w_k \mid t_k,\, N_e(\cdot)\right)
    = \frac{\binom{k}{2}}{N_e(t_k + w_k)}\,
      \exp\!\left(-\binom{k}{2}\int_{t_k}^{t_k + w_k}
      \frac{dt}{N_e(t)}\right),
    \label{eq:varcoal}
\end{equation}

and the full log-likelihood is obtained by summing the logarithms of
Eq.~\eqref{eq:varcoal} over all $n-1$ coalescent intervals.

\medskip
\noindent\textbf{Piecewise-constant approximation.}
A tractable and widely used model represents $N_e(t)$ as a piecewise constant function
that changes values only at a set of pre-specified grid points
$0 = x_0 < x_1 < \cdots < x_M$ \cite{gill2013}:

\begin{equation}
    N_e(t) = \theta_j \quad \text{for } t \in [x_{j-1},\, x_j),
    \quad j = 1, \ldots, M+1.
    \label{eq:piecewise}
\end{equation}

Within each grid interval, the integral in Eq.~\eqref{eq:varcoal} reduces to a sum of
terms of the form $\binom{k}{2}\,\Delta t / \theta_j$, recovering the constant-$N_e$
structure and enabling efficient recursive computation during tree traversal.

The joint log-likelihood over all intervals and all $n-1$ coalescent events is:

\begin{equation}
    \ell\!\left(\boldsymbol{\theta} \mid \mathcal{T}\right)
    = \sum_{j=1}^{M+1}\left[-c_j\log\theta_j - \frac{\mathrm{SS}_j}{\theta_j}\right]
    + \text{const},
    \label{eq:pccoal}
\end{equation}

where $c_j$ is the number of coalescent events in interval $j$, and
$\mathrm{SS}_j = \sum_k \binom{k}{2}\,\Delta t_{k,j}$ accumulates the
coalescent-pair--weighted duration of all lineage intervals that overlap grid cell $j$.

\subsubsection{Connection to Birth-Death Models}

The birth-death model of Nee et al.\ (1994) \cite{nee1994} provides an alternative
likelihood framework for reconstructed phylogenies, parameterised by a per-lineage
speciation (birth) rate $\lambda$ and extinction (death) rate $\mu$.
Under complete sampling ($\rho = 1$) of a clade of $n$ extant taxa, the likelihood
of the reconstructed tree with internal node ages
$x_1 \geq x_2 \geq \cdots \geq x_{n-1}$ (measured from the present) is

\begin{equation}
    \mathcal{L}(\lambda, \mu \mid \mathcal{T})
    = (n-1)!\;\lambda^{n-2}
      \prod_{i=1}^{n-1} \Phi(x_i),
    \label{eq:bdlik}
\end{equation}

where the branching-time density is

\begin{equation}
    \Phi(t) = \frac{(\lambda - \mu)^2\,e^{(\lambda-\mu)t}}
               {\left[\lambda\,e^{(\lambda-\mu)t} - \mu\right]^2}.
    \label{eq:phi}
\end{equation}

For incomplete random sampling with fraction $\rho \in (0,1]$, the probability
$p_1(t)$ that a single lineage alive at time $t$ in the past has at least one
sampled descendant today is

\begin{equation}
    p_1(t) = \frac{\rho\,(\lambda - \mu)}
              {\rho\lambda + \bigl[\lambda(1-\rho) - \mu\bigr]\,e^{-(\lambda-\mu)t}},
    \label{eq:p1}
\end{equation}

and Eq.~\eqref{eq:bdlik} is modified accordingly.

\medskip
\noindent\textbf{Epidemiological interpretation.}
In the context of infectious disease, $\lambda$ represents the per-lineage transmission
rate and $\mu$ the per-lineage removal (recovery or death) rate, so that
$R_0 = \lambda/\mu$ is the basic reproduction number. The sampling fraction $\rho$
represents the proportion of infected individuals whose sequences were obtained.
Maximising $\ell(\lambda, \mu, \rho \mid \mathcal{T})$ over the observed phylogeny
of pathogen sequences thus directly estimates key epidemiological parameters.

\subsubsection{Benchmarking Framework for Likelihood Evaluation}

The following computational steps are required to evaluate the coalescent (or
birth-death) log-likelihood given a phylogenetic tree $\mathcal{T}$:

\begin{enumerate}
    \item \textbf{Tree traversal} (postorder): collect the $n - 1$ internal node
          heights $x_i$ and compute coalescent waiting times $w_k$ along the way.
          This step drives the performance cost and motivates Part~I.
    \item \textbf{Sufficient statistic computation}: accumulate $S$ (Eq.~\ref{eq:sstat})
          or, for variable $N_e$, the per-interval sums $\mathrm{SS}_j$
          (Eq.~\ref{eq:pccoal}).
    \item \textbf{Likelihood evaluation}: compute Eq.~\eqref{eq:loglik} or
          Eq.~\eqref{eq:pccoal} for the coalescent model, or Eqs.~\eqref{eq:bdlik}--\eqref{eq:phi}
          for the birth-death model.
    \item \textbf{Optimisation}: maximise the log-likelihood with respect to $N_e$
          (closed-form via Eq.~\ref{eq:MLE}) or with respect to $(\lambda, \mu, \rho)$
          numerically.
\end{enumerate}

Benchmark specifications:
\begin{itemize}
    \item Constant-$N_e$ coalescent likelihood evaluation on the 1{,}000-tip ultrametric
          trees generated in Part~I.
    \item Variable-$N_e$ piecewise-constant likelihood (Skyline model) with $M = 10$
          and $M = 50$ grid points.
    \item Birth-death likelihood evaluation on the same tree sets.
    \item Comparison of \ape (\texttt{birthdeath()}), \castor, and \dendropy
          implementations.
    \item Assessment of accuracy of $\hat{N}_e$ and $(\hat\lambda, \hat\mu)$ under
          complete and incomplete sampling regimes.
\end{itemize}

% ==========================================
% 3. RESULTS AND DISCUSSION
% ==========================================
\section{Results and Discussion}

% ------------------------------------------
\subsection{Tree Traversal Performance}
% ------------------------------------------

\subsubsection{Overview}

Benchmarking results across eight tools reveal substantial performance variation,
spanning from $\sim 50\,\mu\text{s}$ (\cogent preorder on 1k-tree) to
$\sim 70\,\text{ms}$ (\ete ultrametric check on 1k-tree).
A summary comparison is provided in Table~\ref{tab:summary}.

\begin{table}[h!]
\centering
\caption{Summary of tool performance rankings across operations.
         Median times in microseconds (1k-leaf ultrametric trees).}
\label{tab:summary}
\begin{tabular}{lcccc}
\toprule
\textbf{Tool} & \textbf{Preorder} & \textbf{Postorder} & \textbf{Level-order} & \textbf{Ultrametric} \\
\midrule
\cogent       & 148     & 404     & 161     & 261     \\
\phylo        & 289     & 3{,}023 & 1{,}296 & 3{,}258 \\
\toytree      & 2{,}086 & 1{,}858 & 1{,}790 & 5{,}756 \\
\ete          & 2{,}196 & 7{,}921 & 1{,}998 & 70{,}056 \\
\dendropy     & 4{,}466 & 3{,}836 & 3{,}790 & 4{,}954 \\
\biopython    & 28{,}054 & 28{,}477 & 7{,}360 & 3{,}526 \\
\bottomrule
\end{tabular}
\end{table}

\subsubsection{Compiled vs.\ Interpreted Implementations}

\paragraph{High-Performance Tools (\cogent, \toytree)}

\cogent exhibits exceptional performance for traversal operations
($\sim 150$--$400\,\mu\text{s}$), likely due to algorithmic optimisation or compiled
extensions in critical sections. \toytree also demonstrates strong performance
($\sim 1.8$--$2.1\,\mu\text{s}$ for traversal, $\sim 5.8\,\mu\text{s}$ for
ultrametric check), consistent with its Cython-based implementation.

\paragraph{Pure Python Tools (\dendropy, \biopython, \ete)}

\dendropy achieves moderate performance ($\sim 3.8$--$4.5\,\mu\text{s}$), acceptable
for interactive workflows but substantially slower than compiled alternatives.
\biopython exhibits variable performance: excellent on ultrametricity checking
($\sim 3.5\,\mu\text{s}$) but slow on traversals ($\sim 28\,\mu\text{s}$).
\ete shows dramatic overhead on ultrametricity checks ($\sim 70\,\mu\text{s}$),
suggesting algorithmic differences or missing optimisations.

\paragraph{Julia (\phylo)}

\phylo demonstrates moderate traversal performance ($\sim 289\,\mu\text{s}$
preorder) with high variance, reflecting JIT compilation warm-up costs.
Postorder traversal is notably slower ($\sim 3\,\mu\text{s}$), indicating potential
algorithmic inefficiency. The high standard deviation ($\sim 5{,}100\,\mu\text{s}$)
makes performance predictions unreliable on first-run executions.

\subsubsection{Tree Size Scaling}

Evaluation on larger trees (10{,}000 leaves) to assess algorithmic complexity and
scalability is pending. Expected asymptotic complexity is $O(n)$ for all traversal
methods; deviations would indicate implementation inefficiencies (e.g.\ quadratic
list operations).

% ------------------------------------------
\subsection{Coalescent Likelihood Evaluation}
% ------------------------------------------

\subsubsection{Constant-$N_e$ MLE Accuracy}

The closed-form estimator $\hat{N}_e$ (Eq.~\ref{eq:MLE}) was evaluated on the 1{,}000
ultrametric trees generated with \texttt{rcoal(1000)} under a known true $N_e = 1$.
Since all trees were generated under Kingman's coalescent, $\hat{N}_e$ is expected to
be approximately unbiased.
Full results (bias, variance, coverage of confidence intervals) to be reported.

\subsubsection{Variable-$N_e$ Piecewise Likelihood}

Implementation of the Skyline log-likelihood (Eq.~\ref{eq:pccoal}) across tools,
evaluating both correctness (agreement of log-likelihood values) and computational
cost as a function of the number of grid points $M$.
To be reported.

\subsubsection{Birth-Death Parameter Estimation}

Evaluation of the birth-death likelihood (Eqs.~\ref{eq:bdlik}--\ref{eq:p1}) on
heterochronic trees simulated under known $(\lambda, \mu, \rho)$ parameters.
Performance and accuracy comparison between \ape \texttt{birthdeath()}, \castor, and a
pure-Python reference implementation.
To be reported.

\subsubsection{Sampling Bias Effects}

Assessment of parameter estimation bias under non-random (e.g.\ clustered) sampling
regimes, where the incomplete-sampling correction (Eq.~\ref{eq:p1}) is particularly
important.
To be reported.

% ==========================================
% 4. CONCLUSIONS AND PERSPECTIVES
% ==========================================
\section{Conclusions and Perspectives}

\subsection{Key Findings}

\begin{enumerate}
    \item \textbf{Substantial Performance Variation:} Performance differences exceed
          400-fold across tools, emphasising the importance of tool selection for
          computationally intensive phylogenetic analyses.

    \item \textbf{Compiled Backends Essential:} Tools with compiled implementations
          (\cogent, \toytree, \phylo JIT) outperform pure Python alternatives by
          $5$--$200\times$.

    \item \textbf{Algorithm Choice Matters:} Identical operations show significant
          variance across implementations (e.g.\ \ete's ultrametricity check).

    \item \textbf{Language Trade-offs:} R-based tools (\ape) balance performance with
          extensive statistical capabilities. Julia shows promise but current
          implementations exhibit high variance, likely reflecting JIT overhead.

    \item \textbf{Coalescent Likelihood:} The closed-form structure of the constant-$N_e$
          likelihood (Eq.~\ref{eq:MLE}) reduces inference to a single postorder
          traversal, making the choice of traversal backend a direct determinant of
          likelihood computation throughput.
\end{enumerate}

\subsection{Recommendations}

\begin{itemize}
    \item \textbf{Interactive analysis:} \ete or \toytree for visualisation and usability.
    \item \textbf{Performance-critical traversal:} \cogent or \toytree for maximum
          throughput.
    \item \textbf{Coalescent/birth-death likelihood:} \ape (with C backend) or \castor
          for large-scale inference.
    \item \textbf{Future consideration:} Investigate Rust-based implementations
          (e.g.\ \phylors) for extreme-scale datasets.
\end{itemize}

\subsection{Future Directions}

\begin{enumerate}
    \item Expand benchmarking to tools in other languages (Rust: phylo-rs, Go: gotree, etc.). 
    \item Complete benchmarking of coalescent and birth-death likelihood evaluation
          across tools.
    \item Evaluate performance on real outbreak datasets.
    \item Investigate memory usage and scalability limits.
    % items added by pasteur ai (verify pertinence/relevance)
    % \item Assess accuracy of $(\hat\lambda, \hat\mu, \hat\rho)$ under non-random
          sampling regimes.
   %  \item Profile code hotspots to identify optimisation opportunities in the
          likelihood loop.
    \item Explore hybrid approaches combining tools' strengths (e.g.\ Julia for
          inference, Python for visualisation).
\end{enumerate}
! item 2 should be done by end of internship 

% ==========================================
% REFERENCES
% ==========================================
\section{References}

\begin{thebibliography}{99}

\bibitem{paradis2004ape}
Paradis, E., Claude, J., \& Strimmer, K. (2004).
APE: analyses of phylogenetics and evolution in R language.
\textit{Bioinformatics}, 20(2), 289--290.

\bibitem{sukumaran2010dendropy}
Sukumaran, J., \& Holder, M.~T. (2010).
DendroPy: a Python library for phylogenetic computing.
\textit{Bioinformatics}, 26(12), 1569--1571.

\bibitem{knight2018cogent3}
Mucaki, E.~J., \& others (2018).
Cogent3: a toolkit for biological sequence analysis.
\textit{Molecular Biology and Evolution}.

\bibitem{cock2009biopython}
Cock, P.~J.~A., et al.\ (2009).
Biopython: freely available Python tools for computational biology and
bioinformatics.
\textit{Bioinformatics}, 25(3), 422--423.

\bibitem{huerta2016ete3}
Huerta-Cepas, J., Serra, F., \& Bork, P. (2016).
ETE 3: reconstruction, analysis, and visualization of phylogenomic data.
\textit{Molecular Biology and Evolution}, 33(6), 1635--1638.

\bibitem{eaton2020toytree}
Eaton, D.~A. (2020).
toytree: A lightweight Python library for phylogenetics.
\textit{bioRxiv}.

\bibitem{kingman1982}
Kingman, J.~F.~C. (1982).
On the genealogy of large populations.
\textit{Journal of Applied Probability}, 19(A), 27--43.

\bibitem{griffiths1994}
Griffiths, R.~C., \& Tavar\'{e}, S. (1994).
Sampling theory for neutral alleles in a varying environment.
\textit{Philosophical Transactions of the Royal Society B}, 344, 403--410.

\bibitem{nee1994}
Nee, S., May, R.~M., \& Harvey, P.~H. (1994).
The reconstructed evolutionary process.
\textit{Philosophical Transactions of the Royal Society B}, 344, 305--311.

\bibitem{gill2013}
Gill, M.~S., et al.\ (2013).
Improving Bayesian population dynamics inference: a coalescent-based model
for multiple loci.
\textit{Molecular Biology and Evolution}, 30(3), 713--724.

\bibitem{bezanson2017julia}
Bezanson, J., Edelman, A., Karpinski, S., \& Shah, V.~B. (2017).
Julia: A fresh approach to numerical computing.
\textit{SIAM Review}, 59(1), 65--98.
fin
\bibitem{teufel2018castor}
Louca, S., \& Doebeli, M. (2018).
Efficient comparative phylogenetics on large trees.
\textit{Bioinformatics}, 34(6), 1053--1055.

\bibitem{volz2009}
Volz, E. M., et al. (2009).
Viral phylodynamics.
\textit{PLoS Computational Biology}, 5(3), e1000809.

\bibitem{pybus2001}
Pybus, O. G., et al. (2001).
Inferring the population size and growth history of HIV-1 from gene genealogies.
\textit{Journal of Virology}, 75(16), 7529--7539.

\end{thebibliography}

\newpage

% ==========================================
% APPENDIX
% ==========================================
\section{Appendix}

\appendix

\section{Specifications}

\subsection{System Specifications}
\begin{itemize}
    \item OS: Linux Mint 22.3 x86\_64, 6.8.0-111-generic 
    \item Host: OMEN Gaming Laptop 16-ap0xxx 
    \item CPU: AMD Ryzen 9 8940HX with Radeon Graphics (32) @ 5.385GHz 
    \item GPU: AMD ATI 07:00.0 Raphael, NVIDIA 01:00.0 NVIDIA Corporation Device 2d19 
    \item RAM: 32 GB
\end{itemize}

\subsection{Tool Specifications}

\begin{itemize}
    \item \ape: Version 5.8.1; R 4.x, C backend
    \item \castor: Version 1.8.5; R 4.x, C backend
    \item \dendropy: Version 5.0.8; Python 3.x, pure Python
    \item \cogent: Version 2026.5.19a0; Python 3.x, optimised/compiled sections
    \item \biopython: Version 1.87; Python 3.x, pure Python with wrappers
    \item \ete: Version 3.1.3; Python 3.x, pure Python + compiled extensions (version-dependent)
    \item \toytree: Version 3.0.11; Python 3.x, Cython + C backend
    \item \phylo: Version 0.5.4; Julia 1.x, LLVM JIT compilation
\end{itemize}

\subsection{Derivation of the Constant-$N_e$ MLE}

From Eq.~\eqref{eq:loglik}, differentiating with respect to $N_e$:

\begin{equation*}
    \frac{\partial \ell}{\partial N_e}
    = -\frac{n-1}{N_e} + \frac{S}{N_e^2} = 0
    \implies
    \hat{N}_e = \frac{S}{n-1}.
\end{equation*}

The second derivative $\partial^2\ell/\partial N_e^2 = (n-1)/N_e^2 - 2S/N_e^3$
evaluated at $\hat{N}_e$ gives $-(n-1)/\hat{N}_e^2 < 0$, confirming a maximum.
The Fisher information is $\mathcal{I}(N_e) = (n-1)/N_e^2$, so the asymptotic
variance of $\hat{N}_e$ is $N_e^2/(n-1)$.

\section{Detailed Results}

\subsection{Performance Plots}

! Violin plots showing traversal time distributions for each tool to be
inserted here. Files: \texttt{ape\_traverse\_ultra.png},
\texttt{*\_traverse\_hetero.png},
\texttt{*\_traverse\_ultra.png},

\subsection{Raw Benchmark Data Tables}

! Detailed tables with all timing measurements, statistical summaries, and system
specifications to be included.

\end{document}
